% Bayes-rule schematic: a 2x2 contingency table for a diagnostic test
% with sensitivity, specificity, prior, and posterior, populated for
% a 1% prior, 95% sensitivity, 98% specificity case to match the
% running example in section 13.2.
\begin{figure}[htbp]
\centering
\begin{tikzpicture}[
  cell/.style={draw=engblack, minimum height=0.95cm,
    minimum width=2.4cm, align=center, font=\small},
  header/.style={draw=engblack, minimum height=0.7cm,
    minimum width=2.4cm, align=center, font=\small\bfseries,
    fill=engblue!10},
  arrow/.style={-{Stealth}, thick, draw=engblack!70},
]

% Title row (population is 100 000 to keep counts whole)
\node[anchor=south, font=\small\bfseries] at (3.6, 4.4)
  {Population of $100\,000$ with prior $\Pr(D) = 0.01$};
\node[anchor=south, font=\scriptsize] at (3.6, 4.0)
  {sensitivity $\Pr(+ \mid D) = 0.95$, specificity $\Pr(- \mid D^{c}) = 0.98$};

% Column headers
\node[header] (h1) at (2.4, 3.0) {$D$ (disease)};
\node[header] (h2) at (4.8, 3.0) {$D^{c}$ (no disease)};

% Row headers
\node[header, minimum width=1.6cm] (r1) at (0.6, 2.05) {$+$ test};
\node[header, minimum width=1.6cm] (r2) at (0.6, 1.05) {$-$ test};

% Cells
\node[cell, fill=engred!10] (c11) at (2.4, 2.05)
  {true positives\\$1{,}000 \cdot 0.95 = 950$};
\node[cell, fill=engamber!12] (c12) at (4.8, 2.05)
  {false positives\\$99{,}000 \cdot 0.02 = 1{,}980$};
\node[cell, fill=engamber!12] (c21) at (2.4, 1.05)
  {false negatives\\$1{,}000 \cdot 0.05 = 50$};
\node[cell, fill=engblue!10] (c22) at (4.8, 1.05)
  {true negatives\\$99{,}000 \cdot 0.98 = 97{,}020$};

% Marginal totals
\node[cell, minimum width=2.4cm, fill=engblue!5] (tot1) at (7.4, 2.05)
  {total $+$:\\$2{,}930$};
\node[cell, minimum width=2.4cm, fill=engblue!5] (tot2) at (7.4, 1.05)
  {total $-$:\\$97{,}070$};

\draw[arrow] (c12.east) -- (tot1.west);
\draw[arrow] (c22.east) -- (tot2.west);

% Posterior calculation
\node[cell, minimum width=10cm, minimum height=1.2cm,
  fill=englime!15] at (4.4, -0.3)
  {Posterior $\Pr(D \mid +) = \dfrac{950}{2{,}930} \approx 0.324$.
   A positive test moves belief from $1\,\%$ to about $32\,\%$.};

\end{tikzpicture}
\caption[Bayes-rule schematic on a diagnostic test]{The Bayes-rule
calculation for a screening test with prior $\Pr(D) = 0.01$,
sensitivity $0.95$, and specificity $0.98$. The four cells of the
contingency table are obtained by multiplying the prior cell sizes
by the test characteristics. The posterior is the ratio of true
positives to total positives; the false-positive count dominates
the denominator because the disease-free population is much larger
than the diseased one, even with high specificity. This is the
mechanism behind the base-rate fallacy named in
section~\ref{sec:vol02:ch13:bayes}. Source: computed by the
authors; see also \cite[ch.~3]{text:ross-probability}.}
\label{fig:vol02:ch13:bayes-schematic}
\end{figure}
